\documentclass[11pt]{article}
\usepackage{amsmath,amssymb,amsthm,booktabs}
\usepackage[margin=1in]{geometry}
\usepackage[colorlinks=true,linkcolor=blue,citecolor=blue,urlcolor=blue]{hyperref}
\newtheorem{theorem}{Theorem}[section]
\newtheorem{proposition}[theorem]{Proposition}
\newtheorem{lemma}[theorem]{Lemma}
\newcommand{\F}{\mathbb F}
\newcommand{\one}{\mathbf 1}
\newcommand{\Tr}{\mathrm{Tr}}
\newcommand{\Res}{\mathrm{Res}}
\newcommand{\abs}[1]{\lvert#1\rvert}
\newcommand{\Lstar}{\tfrac{\sqrt e}{2}}
\title{Sharp determinant bounds for three\\ near-Hadamard families}
\author{Jeffery Kline}
\date{Version 0.1.0 --- release candidate, November 15, 2026}
\begin{document}
\maketitle

\begin{abstract}
No $n\times n$ matrix with entries $\pm1$, where $n\equiv2\pmod4$, has
determinant larger in absolute value than the Ehlich--Wojtas bound
$W_n=(2n-2)(n-2)^{(n-2)/2}$. For each prime power $q\equiv1\pmod4$,
Momihara, Suda and Xiang constructed three such matrices of order
$n=q+1$, their Classes (4), (5) and (6). They wrote the determinants as
products over Gauss and Jacobi sums without evaluating them, and asked
whether each family reaches at least $0.8\,W_n$ for every $n\ge30$.
A 2026 survey lists this as open. We prove that every matrix of order
$n\ge30$ in the three families exceeds $0.8\,W_n$, and that in each family the ratio to
$W_n$ tends to $\sqrt e/2\approx0.824$, the limit they suggested. We also
find the smallest ratio in each family exactly: about $0.8169$, $0.8154$
and $0.8184$, attained only at $n=30$, $38$ and $50$.

The proof needs only two exact facts about the factors of each product,
their sum and their sum of squares, both given by character
orthogonality, together with the size of a Gauss sum. This settles every
$q\ge676$ for Class (4) and every $q\ge1001$ for Classes (5) and (6).
Exact integer computations settle the remaining 306 determinants. The same
estimate exceeds the three smallest ratios for every $q\ge10525$, and a
rigorous interval computation covers the orders below. We have
not read the full text of a 2026 follow-up by Li, Momihara and Xiang, so
we cannot confirm that the result is new.
\end{abstract}

\section{Introduction}\label{sec:intro}

\subsection{The question}
A square matrix with entries $\pm1$ has the largest possible determinant
when its rows are orthogonal; such a matrix is a Hadamard matrix, and it
can exist only when the order $n$ is $1$, $2$ or a multiple of $4$. For
$n\equiv2\pmod4$, Ehlich and Wojtas showed that
\begin{equation}\label{eq:EW}
 \abs{\det M}\le W_n=(2n-2)(n-2)^{(n-2)/2},
\end{equation}
and the bound is rarely attained. The ratio $\abs{\det M}/W_n$ therefore
measures how good a construction is.

Momihara, Suda and Xiang~\cite{MSX}, abbreviated MSX, built nine
families of near-Hadamard matrices: matrices of order $n\equiv2\pmod4$
whose Gram matrix $MM^{\mathsf T}$ is block diagonal, with diagonal
entries $n$ and off-diagonal entries $\pm2$ in each block. For six of the families they computed the determinant in closed
form. For Classes (4), (5) and (6), of order $n=q+1$ with $q\equiv1\pmod4$
a prime power, they could only write it as a product over Gauss and Jacobi
sums. Their computations gave ratios $0.844$, $0.846$ and $0.823$ at
$n=30$, and at least $0.8$ for $38\le n\le1000$; they remarked that the
ratio seems to approach $\sqrt e/2$. Their Problem~7.1 asks for sharp lower
bounds and, in particular, whether the ratio is at least $0.8$ for every
$n\ge30$ in each class. The 2026 survey of Egan, Kotsireas, London and
\'O~Cath\'ain~\cite[\S5.6, Problem 15]{EKLO} repeats the $0.8$ question.

\subsection{The result}
\begin{theorem}\label{thm:main}
Let $q\ge29$ be a prime power with $q\equiv1\pmod4$, and let
$M^{(i)}_{q+1}$ be a matrix of Class $i\in\{4,5,6\}$ of
\cite[Theorem 1.6]{MSX}, built from any primitive element. Then
\begin{equation}\label{eq:target}
 R_i(q):=\frac{\abs{\det M^{(i)}_{q+1}}}{W_{q+1}}
 =\frac{\abs{\det M^{(i)}_{q+1}}}{2q(q-1)^{(q-1)/2}}>\frac45 .
\end{equation}
Moreover $R_i(q)\to\sqrt e/2$ as $q\to\infty$, for each $i$ and
uniformly in the primitive element. Finally, over all $q\ge29$ and all
primitive elements, the smallest values are
\begin{equation}\label{eq:infima}
 \min R_4=R_4(29)\approx0.81690,\qquad
 \min R_5=R_5(37)\approx0.81541,\qquad
 \min R_6=R_6(49)\approx0.81843,
\end{equation}
with the exact rational values given in Section~\ref{sec:finite}. Every
other $q\ge29$ gives a strictly larger ratio in the same class.
\end{theorem}

In Class (4) at $q=29$ the minimum comes from one of the two possible
determinants described in Section~\ref{sec:choices}; the other is about
$0.8440$, the value MSX report. The theorem answers Problem~15 of the
survey, which is the $0.8$ question of Problem~7.1. It also gives sharp
constants in the sense of Problem~7.1: \eqref{eq:infima} gives, for each
class, the largest constant $c_i$ with $\abs{\det M^{(i)}_n}\ge c_iW_n$
for all $n\ge30$, and the limit $\sqrt e/2$ is the best constant for
large $n$.

\subsection{How the proof works}
In each class, $\abs{\det M}$ is a product of numbers $q+1-T$, one for
each character of a cyclic group, where the $T$ are real character sums.
Dividing by $q+1$ gives a product of factors $1-t$ with $t=T/(q+1)$.
Three facts about the $T$ suffice:
\begin{itemize}
\item each satisfies $\abs T\le2\sqrt q$, from the size of a Gauss sum;
\item their sum is a small explicit constant;
\item their sum of squares is close to $q^2$.
\end{itemize}
The second and third facts are exact, and both follow from character
orthogonality because the underlying $\pm1$ sequence is symmetric under
inversion and vanishes only at the identity. Expanding
$\log\prod(1-t)$ to second order then gives $\log R\approx 1-\log2-\tfrac12$,
so $R\approx\sqrt e/2$, with an error of order $q^{-1/2}$. No information
about the individual phases of the Gauss or Jacobi sums is needed.
Explicit constants make the error small enough for $q\ge676$ in Class (4)
and $q\ge1001$ in Classes (5)--(6). Below these thresholds, exact integer
resultants evaluate all 306 determinants that remain. For the minima the
same estimate is used with the target $0.8$ replaced by the smallest
computed ratio; it then needs $q$ up to about $10^4$, and the orders below
are checked in interval arithmetic, which bounds the error of every
operation.

\subsection{Earlier work and status}
The constructions and the determinant products are those of
MSX~\cite[Theorem 1.6, Propositions 2.5, 2.6, 2.8, 3.5 and 4.2]{MSX}.
MSX suggested the limit $\sqrt e/2$ for Classes (4)--(6) and proved the
same limit for their Class (9). Bounding a determinant through the first
two moments of its spectrum is standard: MSX use a two-moment bound in
their Lemma~2.2 and Proposition~2.1, and Ipsen and Lee~\cite[Theorem 2.6]{IL}
give general trace--log expansions with error bounds. What this paper adds
is the deduction of the two moment identities for these three families,
explicit thresholds, the exact computation that closes the finite
range, and the exact minima.

\textbf{Novelty is unconfirmed.} The survey, revised on 3 November 2026,
still lists the question as open while citing the 2026 article of Li,
Momihara and Xiang~\cite{LMX}. We have read only the public preview of that
article: its abstract, introduction and section openings concern type I
near-Hadamard matrices, whose ratio tends to $1$, and do not mention
Classes (4)--(6) or Problem~7.1. We have not read its full text and found
no earlier proof elsewhere, but neither fact rules out a prior solution.

\section{A logarithmic estimate}\label{sec:lemma}

All logarithms are natural. The following is the scalar form of the
trace--log argument, with its error made explicit.

\begin{lemma}\label{lem:log}
Let $t_1,\ldots,t_N$ be real with $\abs{t_j}\le\rho<1$, and put
$S=\sum_jt_j$ and $V=\sum_jt_j^2$. Then
\begin{equation}\label{eq:logerror}
 \Bigl|\sum_{j=1}^N\log(1-t_j)+S+\frac V2\Bigr|
 \le\frac{\rho V}{3(1-\rho)}.
\end{equation}
\end{lemma}
\begin{proof}
For $\abs t\le\rho$ the series of $\log(1-t)$ after its quadratic term is
bounded by $\sum_{k\ge3}\abs t^k/k\le\abs t^3/(3(1-\rho))$. Sum over $j$
and use $\sum_j\abs{t_j}^3\le\rho V$.
\end{proof}

In every application below
\begin{equation}\label{eq:rho}
 \rho=\frac{2\sqrt q}{q+1}<1,
 \qquad \frac{\rho}{3(1-\rho)}=\frac{2\sqrt q}{3(\sqrt q-1)^2},
\end{equation}
and $V<1$, so the error in~\eqref{eq:logerror} is at most
$2\sqrt q/(3(\sqrt q-1)^2)$.

\section{Classes (5) and (6)}\label{sec:56}

Fix a primitive element $w$ of $\F_q$, put $m=(q-1)/2$, and let
$H=\langle w^2\rangle$, the group of nonzero squares. Let $\eta$ be the
quadratic character of $\F_q$ and let $\chi$ be the multiplicative
character with $\chi(w)=e^{2\pi i/(q-1)}$. In this section every
multiplicative character, including the trivial one, is extended by zero
at $0$. Define the Jacobi sums
\[
 J(a,b)=\sum_{x\in\F_q}a(x)\,b(1-x),\qquad
 U_j=J(\chi^j,\eta),\qquad V_j=J(\eta\chi^j,\eta)\qquad(1\le j<m).
\]
By MSX~\cite[Theorems 1.6 and 4.1 and Propositions 2.6, 2.8 and 4.2]{MSX},
\begin{align}
 \abs{\det M^{(5)}_{q+1}}&=(q-1)^2\prod_{j=1}^{m-1}(q+1+U_j+V_j),
   \label{eq:D5}\\
 \abs{\det M^{(6)}_{q+1}}&=(q^2+2q+5)\prod_{j=1}^{m-1}(q+1-U_j-V_j).
   \label{eq:D6}
\end{align}
MSX extend the trivial character differently at $0$, but every character
in these products is nontrivial, so the formulas are unaffected. We work
from the published products and do not need to rebuild the matrices.

\begin{proposition}\label{prop:mom56}
The numbers $X_j=U_j+V_j$ are real and satisfy
\begin{equation}\label{eq:mom56}
 \abs{X_j}\le2\sqrt q,\qquad
 \sum_{j=1}^{m-1}X_j=2,\qquad
 \sum_{j=1}^{m-1}X_j^2=q^2-4q-1.
\end{equation}
\end{proposition}
\begin{proof}
Since $\eta(-1)=1$, substituting $x=y^{-1}$ in $\overline{U_j}$ gives
\[
 \overline{U_j}=\sum_{y\ne0}\chi^j(y)\,\eta(1-y^{-1})
 =\sum_y\chi^j(y)\,\eta(y)\,\eta(1-y)=V_j ,
\]
so $X_j=2\,\mathrm{Re}\,U_j$ is real. The characters $\chi^j$, $\eta$ and
$\chi^j\eta$ are nontrivial for $1\le j<m$, so the Gauss--Jacobi identity
and the modulus of a nontrivial Gauss sum give $\abs{U_j}=\sqrt q$, and
hence $\abs{X_j}\le2\sqrt q$.

For the moments, let $g(x)=\eta(x-1)$ for $x\in H$ and
\[
 Y_j=\sum_{x\in H}g(x)\,\chi^j(x)\qquad(0\le j<m).
\]
The restrictions of $\chi^j$ to $H$, $0\le j<m$, are all the characters of
$H$. The function $g$ satisfies $g(x^{-1})=g(x)$, $g(1)=0$, and $g(x)^2=1$
for $x\ne1$. Writing the indicator of the squares as $(1+\eta)/2$ gives
$2Y_j=U_j+V_j=X_j$ for $1\le j<m$. With our convention at zero,
$J(1,\eta)=J(\eta,\eta)=-1$, so $Y_0=-1$: the first value is the sum of
$\eta(1-x)$ over $x\ne0$, and for the second the substitution
$u=x/(1-x)$ on $\F_q\setminus\{1\}$ gives
$J(\eta,\eta)=\sum_{u\ne-1}\eta(u)=-\eta(-1)=-1$.

Character orthogonality on $H$ now gives
\[
 \sum_{j=0}^{m-1}Y_j=m\,g(1)=0,\qquad
 \sum_{j=0}^{m-1}Y_j^2=m\sum_{x\in H}g(x)\,g(x^{-1})=m(m-1).
\]
Removing $Y_0=-1$ and using $X_j=2Y_j$ gives $\sum X_j=2$ and
$\sum X_j^2=4m(m-1)-4=q^2-4q-1$.
\end{proof}

Write $s=-1$ for Class (5) and $s=1$ for Class (6), and put
$K_{-1}=(q-1)^2$ and $K_{1}=q^2+2q+5$. Every factor
in~\eqref{eq:D5}--\eqref{eq:D6} is at least $q+1-2\sqrt q=(\sqrt q-1)^2>0$.
Dividing by $W_{q+1}=2q(q-1)^{(q-1)/2}$ and by $(q+1)^{m-1}$ gives
\begin{equation}\label{eq:P56}
 R_{s}(q)=C_s(q)\prod_{j=1}^{m-1}(1-t_j),\qquad
 C_s(q)=\frac{K_s}{2q(q-1)}\Bigl(\frac{q+1}{q-1}\Bigr)^{(q-3)/2},\qquad
 t_j=\frac{sX_j}{q+1},
\end{equation}
with $R_{-1}=R_5$ and $R_1=R_6$. By Proposition~\ref{prop:mom56},
\[
 S=\sum_jt_j=\frac{2s}{q+1},\qquad V=\sum_jt_j^2=\frac{q^2-4q-1}{(q+1)^2}<1.
\]
The inequalities $\log(1-x)\ge-x/(1-x)$ for $0<x<1$ and
$\log\frac{1+x}{1-x}\ge2x$ give
\[
 \log\frac{K_s}{2q(q-1)}\ge-\log2-\frac1{q-1},\qquad
 \frac{q-3}{2}\log\frac{q+1}{q-1}\ge1-\frac3q .
\]
Combining these with Lemma~\ref{lem:log}, $-S\ge-2/(q+1)$ and $V<1$ gives,
for $i=5,6$ and every admissible $q\ge5$,
\begin{equation}\label{eq:E56}
 R_i(q)\ge\Lstar\,e^{-E_{56}(q)},\qquad
 E_{56}(q)=\frac1{q-1}+\frac3q+\frac2{q+1}+\frac{2\sqrt q}{3(\sqrt q-1)^2}.
\end{equation}

For the limit, $C_s(q)\to e/2$, $S\to0$, $V\to1$, and the error
in~\eqref{eq:logerror} tends to $0$. Hence $\log R_s(q)\to\frac12-\log2$,
that is, $R_s(q)\to\sqrt e/2$. The products do not depend on the primitive
element at all: changing $w$ only relabels the characters of $H$.

\section{Class (4)}\label{sec:4}

Put $v=(q+1)/2$, fix a primitive element $\omega$ of $\F_{q^2}$, and define
\[
 G=\langle\omega^{2(q-1)}\rangle,\qquad \alpha=\omega^v,\qquad
 D_i=\bigl\{x\in G:\Tr_{q^2/q}(\alpha^ix)\text{ is a nonzero square in }\F_q\bigr\}
 \quad(i=0,1).
\]
The group $G$ has odd order $v$. MSX~\cite[Theorem 3.4]{MSX} build Class (4)
from these sets. Write $L=D_1$ and index matrices by $G$: let
$A_{x,y}=2\one_{D_0}(yx^{-1})-1$, $B_{x,y}=2\one_L(yx^{-1})-1$ and
$(N_L)_{x,y}=\one_L(yx^{-1})$. By~\cite[Proposition 2.5]{MSX},
\begin{equation}\label{eq:gram4}
 M=\begin{pmatrix}A&B\\B^{\mathsf T}&-A^{\mathsf T}\end{pmatrix},
 \qquad MM^{\mathsf T}=\mathrm{diag}(C,C),
 \qquad C=2J+(q-1)I-4N_L,
\end{equation}
where $J$ is the all-ones matrix.

For $x\in G$ we have $x^q=x^{-1}$, and $\alpha^{q-1}=-1$, so
\[
 \Tr_{q^2/q}(\alpha x)=\alpha(x-x^{-1}).
\]
This vanishes only at $x=1$, since $G$ has odd order. Inversion negates
the trace and $\eta(-1)=1$, so $L=L^{-1}$ and $1\notin L$. Define
\[
 f(x)=2\one_L(x)-1+\one_{\{1\}}(x),\qquad
 T_\psi=2\sum_{x\in G}f(x)\,\psi(x)\qquad(\psi\in\widehat G).
\]
Then $f(1)=0$, $f(x)^2=1$ for $x\ne1$, and $f(x^{-1})=f(x)$, so every
$T_\psi$ is real. Here the trivial character is included.

\begin{proposition}\label{prop:mom4}
For each $\psi\in\widehat G$, the eigenvalue of $C$ on $\psi$ is
$q+1-T_\psi$, and
\begin{equation}\label{eq:mom4}
 \abs{T_\psi}\le2\sqrt q,\qquad
 \sum_{\psi\in\widehat G}T_\psi=0,\qquad
 \sum_{\psi\in\widehat G}T_\psi^2=q^2-1.
\end{equation}
Consequently
\begin{equation}\label{eq:D4}
 \abs{\det M^{(4)}_{q+1}}=\prod_{\psi\in\widehat G}(q+1-T_\psi).
\end{equation}
\end{proposition}
\begin{proof}
Write $\psi(L)=\sum_{x\in L}\psi(x)$. For nontrivial $\psi$,
$T_\psi=4\psi(L)+2$ and the eigenvalue of $C$ is $q-1-4\psi(L)=q+1-T_\psi$.
For the trivial character, $T_1=4\abs L-(q-1)$ and the eigenvalue is
$2q-4\abs L=q+1-T_1$.

For the bound, let $\theta$ be a multiplicative character of $\F_{q^2}$ of
order $v$, and let $\kappa$ have order $4$. Since $\gcd(v,2(q-1))=1$, the
restrictions of $\theta^h$, $0\le h<v$, are all the characters of $G$.
Write $\mathcal G_Q(\cdot)$ for a Gauss sum over $\F_Q$ with the canonical
additive character. MSX~\cite[Proposition 3.5]{MSX} show that
\[
 \theta^h(L)=\frac{\mathcal G_q(\eta)}{4q}
 \sum_{\ell\in\{1,3\}}\mathcal G_{q^2}(\kappa^\ell\theta^h)\,\kappa^{-\ell}(\alpha)
 +\begin{cases}(q-1)/4,&h=0,\\-1/2,&h\ne0.\end{cases}
\]
The constant terms cancel in $T_\psi$. Each $\kappa^\ell\theta^h$ is
nontrivial, because its order-$4$ part cannot cancel a character of odd
order. Nontrivial Gauss sums have modulus $\sqrt q$ over $\F_q$ and $q$
over $\F_{q^2}$, so $\abs{T_\psi}\le(\sqrt q/q)(q+q)=2\sqrt q$.

Character orthogonality gives the moments:
\[
 \sum_\psi T_\psi=2v\,f(1)=0,\qquad
 \sum_\psi T_\psi^2=4v\sum_{x\in G}f(x)\,f(x^{-1})=4v(v-1)=q^2-1.
\]
Every eigenvalue of $C$ is at least $(\sqrt q-1)^2>0$, and
$\abs{\det M}=\det C$ by~\eqref{eq:gram4}, which gives~\eqref{eq:D4}.
\end{proof}

Dividing~\eqref{eq:D4} by $W_{q+1}$ gives
\begin{equation}\label{eq:P4}
 R_4(q)=C_4(q)\prod_{\psi\in\widehat G}(1-t_\psi),\qquad
 C_4(q)=\frac{q+1}{2q}\Bigl(\frac{q+1}{q-1}\Bigr)^{(q-1)/2},\qquad
 t_\psi=\frac{T_\psi}{q+1}.
\end{equation}
Now $S=0$ and $V=(q-1)/(q+1)<1$. The inequalities
$\log(1+x)\ge x-x^2/2$ for $x\ge0$ and $\log\frac{1+x}{1-x}\ge2x$ give
\[
 \log C_4(q)\ge-\log2+\frac1q-\frac1{2q^2}+\frac{q-1}{q}=1-\log2-\frac1{2q^2},
\]
and Lemma~\ref{lem:log} then gives, for every admissible $q\ge5$,
\begin{equation}\label{eq:E4}
 R_4(q)\ge\Lstar\,e^{-E_4(q)},\qquad
 E_4(q)=\frac1{2q^2}+\frac{2\sqrt q}{3(\sqrt q-1)^2}.
\end{equation}
As before, $C_4(q)\to e/2$, $V\to1$ and the error tends to $0$, so
$R_4(q)\to\sqrt e/2$.

\subsection{Choice of primitive element}\label{sec:choices}
Unlike Classes (5)--(6), Class (4) can depend on the primitive element.
Any other primitive element is $\omega^r$ with $\gcd(r,q^2-1)=1$, so $r$ is
odd. The group $G$ is unchanged, and
\[
 \alpha'=\alpha^r=\alpha\,(\omega^{q+1})^{(r-1)/2}.
\]
The extra factor lies in $\F_q^*$ and has quadratic character
$(-1)^{(r-1)/2}$, because $\omega^{q+1}$ generates $\F_q^*$. So $L$ is
replaced either by itself or by $G\setminus(L\cup\{1\})$, that is, $f$ by
$f$ or $-f$, and replacing $\omega$ by $\omega^{-1}$ realizes the second
case. Reordering $G$ only permutes the characters. Each $q$ therefore has
at most two determinant values, one for each sign. The analytic bound
covers both, since~\eqref{eq:mom4} is unchanged by $f\mapsto-f$; the finite
computation checks both.

\section{Thresholds, the finite range and the minima}\label{sec:finite}

\subsection{The analytic thresholds}
Each term of $E_4$ and $E_{56}$ decreases in $q$; for the last term, with
$u=\sqrt q$,
\[
 \frac{d}{du}\,\frac{2u}{3(u-1)^2}=-\frac{2(u+1)}{3(u-1)^3}<0.
\]
At $q=676=26^2$,
\begin{equation}\label{eq:threshold4}
 E_4(676)=\frac1{2\cdot676^2}+\frac{52}{1875}
 =\frac{47527379}{1713660000}<\frac1{35}.
\end{equation}
For $q\ge1001$ we have $\sqrt q>63/2$, so
\begin{equation}\label{eq:threshold56}
 E_{56}(q)\le\frac1{1000}+\frac3{1001}+\frac2{1002}+\frac{84}{3721}
 =\frac{53309553221}{1866085221000}<\frac1{35}.
\end{equation}
Since $e>\sum_{k=0}^5 1/k!=163/60>784/289$, we have $\sqrt e>28/17$, and
with $e^{-x}\ge1-x$ an error below $1/35$ gives
\[
 R_i(q)>\Lstar\cdot\frac{34}{35}=\frac{17\sqrt e}{35}>\frac{17}{35}\cdot\frac{28}{17}=\frac45 .
\]
This proves~\eqref{eq:target} for $q\ge676$ in Class (4) and for
$q\ge1001$ in Classes (5)--(6). The thresholds are not optimized.

\subsection{Exact computation below the thresholds}
For a monic polynomial $h$ with roots $z_1,\ldots,z_d$, the resultant
satisfies $\Res(h,F)=\prod_a F(z_a)$. This turns each determinant product
into an integer resultant.

For Classes (5)--(6), put $g_k=g(w^{2k})=\eta(w^{2k}-1)$ for $0\le k<m$ and
\[
 h(X)=1+X+\cdots+X^{m-1},\qquad
 F_s(X)=q+1-2s\sum_{k=0}^{m-1}g_kX^k .
\]
The roots of $h$ are the values of the nontrivial characters of $H$ at
$w^2$, so $\abs{\det M}=K_s\Res(h,F_s)$ by~\eqref{eq:D5}--\eqref{eq:D6}.
For Class (4), put $\gamma=\omega^{2(q-1)}$, $f_k=f(\gamma^k)$, and
\[
 h_4(X)=X^v-1,\qquad
 F_{4,\epsilon}(X)=q+1-2\epsilon\sum_{k=0}^{v-1}f_kX^k\qquad(\epsilon=\pm1),
\]
so that $\abs{\det M}=\Res(h_4,F_{4,\epsilon})$; here $h_4$ includes the
trivial character, and the two signs are the two cases of
Section~\ref{sec:choices}. Each determinant $\Delta$ passes if the integer
\begin{equation}\label{eq:margin}
 5\Delta-4\cdot2q(q-1)^{(q-1)/2}
\end{equation}
is positive. No floating-point arithmetic enters any decision.

Every prime power $q\equiv1\pmod4$ in the two ranges was checked:
\begin{center}
\begin{tabular}{@{}lrrr@{}}
\toprule
Class & Range & Field orders & Determinants\\
\midrule
(4), both signs & $29\le q<676$ & 64 & 128\\
(5) & $29\le q<1001$ & 89 & 89\\
(6) & $29\le q<1001$ & 89 & 89\\
\bottomrule
\end{tabular}
\end{center}
Every margin~\eqref{eq:margin} is positive. The smallest ratios in these
ranges are
\begin{align*}
 \text{Class (4), }q=29:&\quad
   \tfrac{16067122183125}{19668469112621}\approx0.81690,\\
 \text{Class (5), }q=37:&\quad
   \tfrac{690199704608000}{846441320833557}\approx0.81541,\\
 \text{Class (6), }q=49:&\quad
   \tfrac{13602934789834170625}{16620815899787526144}\approx0.81843.
\end{align*}
Section~\ref{sec:infima} shows that they are the minima over all $q\ge29$.
The Class (4) value at $q=29$ comes from the sign not reported by MSX; the
other sign gives about $0.8440$, matching their value.

Each field $\F_q$ is modelled as a quotient of $\F_p[X]$ by a monic
polynomial. The computation checks that the powers of the chosen element
run through all $q-1$ nonzero residues before returning to $1$, which
proves that the quotient is a field and the element is primitive. For
Class (4), $\F_{q^2}=\F_q[u]/(u^2-d)$ with a certified nonsquare $d$, and
$\omega$ is certified primitive by exact order tests. Field isomorphisms
preserve traces and squares, so one model per order suffices.

A second program, written separately, rebuilds the sets and polynomial
products for all 306 determinants and compares them with full matrix
determinants at $q=29,49,125,625,673$ in Class (4) (both signs) and at
$q=29,49,125,625,729,997$ in Classes (5)--(6); the largest matrix has
order $998$. As a negative control, one sign of Class (4) at $q=13$ gives
$841/1053<4/5$, outside the theorem's range. Both programs use the FLINT
library for exact resultants and determinants, so the arithmetic library
is shared.

Together with~\eqref{eq:threshold4} and~\eqref{eq:threshold56}, this
proves~\eqref{eq:target} for every $q\ge29$.

\subsection{The minima}\label{sec:infima}
Let $\mu_i$ be the smallest computed ratio of Class $i$ above. Since $E_4$
and $E_{56}$ decrease in $q$ and $e^{-x}\ge1-x$, we have $R_i(q)>\mu_i$ for
every $q\ge Q_i$ as soon as $\Lstar\,(1-E_i(Q_i))>\mu_i$, where
$E_5=E_6=E_{56}$. We check this in exact rational arithmetic, replacing
$\sqrt e$ by a rational lower bound derived from $e>\sum_{k\le20}1/k!$ and
$\sqrt{Q_i}$ in the last term of $E_i$ by a rational lower bound. The
smallest cutoffs that pass are
\[
 Q_4=5720,\qquad Q_5=5039,\qquad Q_6=10525 .
\]
Below them lie $392$, $352$ and $666$ admissible orders $q\ge29$. For
these, the factors of each product are a discrete Fourier transform: in
Classes (5)--(6), $X_j=2Y_j$, where $Y_j$ is the transform of the
sequence $g_k$ of Section~\ref{sec:finite}, and in Class (4), $T_\psi$ is
twice the transform of $f_k$. We compute $\log R_i(q)$ in ball arithmetic
(Arb, $128$-bit precision), whose discrete Fourier transform returns
guaranteed enclosures, and check that the whole ball lies above
$\log\mu_i$. For the 306 determinants computed exactly, the ball
contains the exact value and the comparison is made exactly. Every order
other than the minimizer gives a strictly larger ratio, for both signs in
Class (4). This proves~\eqref{eq:infima} and completes the proof of
Theorem~\ref{thm:main}. The computation takes about a minute on one
processor core; computing an exact integer determinant at $q\approx10^4$
takes about two minutes, measured on the same machine.

\section{Remarks}\label{sec:remarks}

\emph{Scope.} The theorem bounds three families constructed by MSX. It does
not construct new matrices or improve the largest known determinant at any
order. The limit $\sqrt e/2$ was suggested by MSX; the contribution is its
proof, the $0.8$ bound and the exact minima.

\emph{Why the moments suffice.} The ratio depends on the phases of
$m-1$ or $v$ Gauss and Jacobi sums, which satisfy arithmetic relations that
are hard to describe completely. The proof avoids them: the bound in
Lemma~\ref{lem:log} depends on the $t$ only through $\max\abs t$, $\sum t$
and $\sum t^2$, and orthogonality fixes the last two exactly, because the
underlying sequence is symmetric under inversion, equals $\pm1$ away from
the identity and vanishes there.

\subsection*{Role of AI}
AI assisted the derivations, the computer verification, the literature
search and the exposition under the author's direction. Process-separated
AI reviews checked the proofs and certificates; they are not independent
expert peer review, and no independent human expert has examined the
proofs. The release repository records the sources, the reviews and their
dispositions.

\begin{thebibliography}{9}
\bibitem{EKLO}
R. Egan, I. S. Kotsireas, S. London, and P. \'O Cath\'ain,
\emph{Recent progress in the construction of Hadamard matrices},
arXiv:2610.27175v2, revised 3 November 2026, Section 5.6, Problem 15.
\url{https://arxiv.org/abs/2610.27175v2}.

\bibitem{IL}
I. C. F. Ipsen and D. J. Lee,
\emph{Determinant approximations}, arXiv:1105.0437, 2011.
\url{https://arxiv.org/abs/1105.0437}.

\bibitem{LMX}
S. Li, K. Momihara, and Q. Xiang,
\emph{New constructions of square $(\pm1)$-matrices with large determinants},
Linear Algebra Appl. \textbf{752} (2027), 66--106; online August 2026.
\url{https://doi.org/10.1016/j.laa.2026.08.013}.

\bibitem{MSX}
K. Momihara, S. Suda, and Q. Xiang,
\emph{Square $(1,-1)$-matrices with large determinants and near-Hadamard matrices},
Linear Algebra Appl. \textbf{654} (2022), 28--55.
\url{https://doi.org/10.1016/j.laa.2022.08.028}.
\end{thebibliography}
\end{document}
